\documentclass[english, oneside]{sapthesis}
\usepackage[utf8]{inputenc}
\usepackage{babel}
\usepackage[hidelinks]{hyperref}
\usepackage{amssymb}
\usepackage{microtype}
\usepackage{subcaption}
\usepackage{listings}
\usepackage{xcolor}
\usepackage{changepage}
\usepackage{multirow}
\usepackage{float}
\usepackage{comment}
\usepackage{tabularx}
\usepackage{flafter}
\usepackage{algorithm}
\usepackage{algpseudocode}
\usepackage[below]{placeins}

% Corregge il bug nativo di \Call con argomenti complessi o formule
\renewcommand{\Call}[2]{\textproc{#1} (#2)}

\setcounter{secnumdepth}{3} % Attiva la numerazione per le subsubsection
\setcounter{tocdepth}{3}    % Inserisce le subsubsection nell'indice

% Definizione della palette di colori per Python
\definecolor{codegreen}{rgb}{0,0.6,0}
\definecolor{codegray}{rgb}{0.5,0.5,0.5}
\definecolor{codepurple}{rgb}{0.58,0,0.82}
\definecolor{backcolour}{rgb}{0.97,0.97,0.97}

% Configurazione globale di listings
\lstset{
    backgroundcolor=\color{backcolour},   
    commentstyle=\color{codegreen}\itshape,
    keywordstyle=\color{blue}\bfseries,
    numberstyle=\tiny\color{codegray},
    stringstyle=\color{codepurple},
    basicstyle=\ttfamily\footnotesize, 
    breakatwhitespace=false,         
    breaklines=true,                 
    postbreak=\mbox{\textcolor{gray}{$\hookrightarrow$}\space}, 
    captionpos=b,                    
    keepspaces=true,                 
    numbers=left,                    
    numbersep=8pt,                  
    showspaces=false,                
    showstringspaces=false,
    showtabs=false,                  
    tabsize=4,
    frame=single,                    
    rulecolor=\color{codegray!30},
}

% --- Candidate and Thesis Metadata ---
\title{Centralized Scheduling in Satellite Computing:\\ An LP-based Approach}
\author{Claudio Bagini}
\IDnumber{2045337}
\course{Computer Science}
\courseorganizer{Faculty of Information Engineering, Computer Science and Statistics}
\submitdate{2025/2026}
\copyyear{2026}
\advisor{Prof. Emiliano Casalicchio}
\authoremail{claudio.bagini@gmail.com}
\versiondate{\today}
\thesistype{Bachelor's Thesis}
\examdate{October/November 2026}
\examiner{Prof. Emiliano Casalicchio}

\begin{document}

\frontmatter
\maketitle

\dedication{To countless cups of coffee, endless iterations,\\
and the version of me that refused to drop the task.}

\begin{abstract}
The evolution of Low Earth Orbit (LEO) satellite megaconstellations has opened new perspectives for Satellite Computing, enabling the extension of edge computing capabilities directly into space (Orbital Edge Computing). However, the highly dynamic network topology, stringent energy constraints of satellite nodes, and strict deadlines of real-time applications make task orchestration and offloading critical challenges.

This thesis presents the design, implementation, and experimental evaluation of a centralized orchestrator based on Integer Linear Programming (ILP). To mitigate solver computational costs and prevent starvation, the system integrates an event-driven dynamic batching mechanism controlled by Batch Size ($BS$) and Batch Timeout ($BT$) thresholds, along with a coordinated formulation of timing constraints that accounts for queuing delay prior to resolution.

Using discrete-event simulations developed within the SimPy framework, the proposed architecture is evaluated under stress conditions across varying request arrival rates ($\lambda$), comparing \emph{Mapping Time} and \emph{Mapping Energy} optimization strategies. The analysis investigates the systemic bottlenecks that emerge under elevated workloads, examining the interplay between physical compute capacity, network saturation, and temporal deadline constraints.

Finally, the centralized orchestrator is benchmarked against state-of-the-art distributed and heuristic solutions (\emph{OrbitAware}, \emph{DTS-base}, \emph{DTS-Optimal}, and distributed ILP) across heterogeneous workload profiles and operational energy budgets. Through this comprehensive evaluation, the thesis aims to characterize the operational boundaries and fundamental trade-offs between the theoretical global optimality of centralized mathematical scheduling and the responsiveness of decentralized models.
\end{abstract}

\tableofcontents

\mainmatter

% --- CHAPTER INCLUSIONS ---
\include{capitoli/chapter1_introduction}
\include{capitoli/chapter2_state_of_the_art}
\include{capitoli/chapter3_ilp_architecture}
\include{capitoli/chapter4_performance_analysis}
\include{capitoli/chapter5_comparative_benchmark}
\include{capitoli/chapter6_conclusions}

\backmatter
\cleardoublepage
\phantomsection
\addcontentsline{toc}{chapter}{Bibliography}
\bibliographystyle{IEEEtran}
\bibliography{bibliografia}

\cleardoublepage
\phantomsection

\addcontentsline{toc}{chapter}{Acknowledgments}


\begin{acknowledgments}
Siamo finalmente giunti alla conclusione di questo percorso interminabile, ancora non me ne rendo conto ma sono vermante passati 4 anni dal suo inizio.
Mi sembra ieri quando mi sono finalmente diplomato ed ho deciso di intraprendere questo nuovo viaggio. 
Questa scelta ha portato tanti cambiamenti dentro di me, mi ha permesso di crescere, di maturare e di approfondire le mie conoscenze, non è stata sempre una passeggiata ma in qualche modo siamo finalemnte giunti alla fine. O quasi\ldots

\medskip

Adesso però comincia la parte più difficile di tutta la tesi, scrivere i ringraziamenti. Non mi trovo molto a mio agio con questo tipo di cose ma ci tengo a 
ringraziare tutte le persone che mi hanno accompagnato e supportato in questo percorso, senza di loro non sarei arrivato a questo punto.
Di conseguenza chiedo a voi lettori di leggere con pazienza e comprensione queste righe, che nonostante siano scritte con il cuore, non riescono a trasmettere tutto quello che vorrei dire.
Vi chiederei di leggere solo la parte a voi indirizzata ma so già che alcuni di voi vorranno leggere tutto, quindi vi chiedo solo di non farmi sentire più in imbarazzo di quanto già io non sia, grazie.

\medskip

In primis vorrei ringraziare la mia famiglia, miei genitori, mia madre, che mi ha sempre supportato ed incoraggiato, e mio padre, che nonostante svariate divergenza di opinioni, mi ha sempre sostenuto e spronato a fare di meglio. Mi avete dato tanto, quando io vi ho ridato poco o niente indietro, ed è un qualcosa che mi porterò sempre nel cuore, grazie di tutto.
Ti ricordi mamma? Di quando mi aiutavi a ripetere storia? O arte? non è mai piaciuta a nessuno di noi 2 ma nonostante tutto tu mi hai sempre aiutato a ripetere. Il tempo dove tu potessi darmi una mano con lo studio è finito da un bel pò ormai ma non dimenticherò mai tutto quello che hai fatto per me.
E tu papà? Ti ricordi durante il liceo? Ogni primo giorno di ogni anno volevi che stessi da te per potermi fare una foto mentre uscivo di casa, sono 4 anni che ormai non fai più quelle foto. Fanne una adesso, in questo momento di conclusione, per ricordarti di quanto io sia cresciuto e di quanto tu mi abbia aiutato a crescere.

\medskip

Voglio ringraziare i miei nonni per aver vegliato su di me per tutti questi anni anche se probabilmente sono io stesso la causa di qualche capello bianco in più che vi siete ritrovati sulla testa, non è vero nonna? Ora che è finita anche questa parte della mia vita che ne pensi di festeggiare? Magari con una bella porzione di polpette al sugo? Questa volta ti darò una mano a prepararle, promesso.

\medskip

Fratello\ldots secondo quanto mi dicono dovrei includere anche te nei ringraziamenti, non so bene per cosa dovrei ringraziarti ma so che in qualche modo mi hai aiutato a crescere, anche se non ne sono sicurissimo ma vabbè, quello di cui sono sicuro è che da fratello maggiore mi hai pavimentato la strada su molte cose come ben sai e per questo ti ringrazio. C'è ne sono voluti di anni prima che iniziassimo a sopportarci ma non mi dispiace questa sorta di alleanza che abbiamo adesso, vediamo quanto durerà, ma per adesso va bene così.

\medskip




\end{acknowledgments} 

\end{document}